\begin{table}[htbp]
    \centering
    \caption{Rancangan Integrasi Protokol ISAPI (\textit{Outbound})}
    \label{tab:api_isapi}
    \small
    \begin{tabularx}{\textwidth}{|>{\raggedright\arraybackslash}p{6cm}|c|X|}
        \hline
        \textbf{ISAPI Path} & \textbf{Metode} & \textbf{Tujuan Operasi} \\
        \hline
        \texttt{/ISAPI/\allowbreak AccessControl/\allowbreak AcsEvent} & POST & Melakukan pemindaian (polling) terhadap log kejadian terbaru. \\
        \hline
        \texttt{/ISAPI/\allowbreak AccessControl/\allowbreak AcsWorkStatus} & GET & Memantau status kesehatan terminal dan kondisi kunci pintu secara real-time. \\
        \hline
        \texttt{/ISAPI/\allowbreak AccessControl/\allowbreak UserInfo/\allowbreak Record} & POST/PUT & Menambahkan atau memperbarui identitas subjek pada memori terminal. \\
        \hline
        \texttt{/ISAPI/\allowbreak AccessControl/\allowbreak UserInfo/\allowbreak Search} & POST & Menarik daftar ID karyawan yang terdaftar untuk fase Mark. \\
        \hline
        \texttt{/ISAPI/\allowbreak AccessControl/\allowbreak UserInfo/\allowbreak Delete} & PUT & Mencabut hak akses subjek dari perangkat keras. \\
        \hline
        \texttt{/ISAPI/\allowbreak AccessControl/\allowbreak RemoteControl/\allowbreak door/1} & PUT & Melakukan instruksi pembukaan pintu secara jarak jauh (remote open). \\
        \hline
        \texttt{/ISAPI/\allowbreak Intelligent/\allowbreak FDLib/\allowbreak FaceDataRecord} & POST/PUT & Melakukan transmisi data biner (multipart) untuk pembaruan profil wajah pada terminal. \\
        \hline
        \texttt{/ISAPI/\allowbreak .../\allowbreak picture/\allowbreak ...} & GET & Mengunduh aset biner foto kejadian anomali berdasarkan referensi URL dinamis. \\
        \hline
    \end{tabularx}
\end{table}